\section*{Bab \ref{ch:b2:heart} --- Jantung dan Daur Jantung}

\begin{solution}{exo:b2:heart:1}
\omterm{def:b2:blood-circulation:vessels}{Vena} cava $\to$ \omterm{def:b2:heart:anatomy}{serambi} kanan $\to$ katup trikuspid $\to$ \omterm{def:b2:heart:anatomy}{bilik}
kanan $\to$ katup paru $\to$ \omterm{def:b2:blood-circulation:vessels}{arteri} paru $\to$ paru-paru
$\to$ \omterm{def:b2:blood-circulation:vessels}{vena} paru $\to$ \omterm{def:b2:heart:anatomy}{serambi} kiri $\to$ katup mitral $\to$ \omterm{def:b2:heart:anatomy}{bilik}
kiri $\to$ katup aorta $\to$ aorta.
\end{solution}

\begin{solution}{exo:b2:heart:2}
Pengisian (diastol): katup AV membuka, semilunarnya menutup. Kerutan
isovolumetrik: keempatnya menutup. Semburan: semilunarnya membuka,
AV-nya menutup. Pengenduran isovolumetrik: semuanya menutup. Bunyi
pertama: katup AV yang menutup pada awal sistolnya; kedua: katup
semilunar yang menutup pada akhirnya.
\end{solution}

\begin{solution}{exo:b2:heart:3}
\omterm{prop:b2:heart:conduction}{Nodus sinoatrium} ($t = 0$; gelombang P ketika \omterm{def:b2:heart:anatomy}{serambinya}
berdepolarisasi); \omterm{prop:b2:heart:conduction}{nodus atrioventrikel} (dicapai pada kira-kira
\qty{80}{ms}, menahan impulsnya selama \qty{100}{ms}: segmen PR-nya);
berkas His dan \omterm{prop:b2:heart:conduction}{serabut Purkinje} (\qtyrange{180}{220}{ms}; QRS ketika
\omterm{def:b2:heart:anatomy}{biliknya} berdepolarisasi); gelombang T menyusul ketika \omterm{def:b2:heart:anatomy}{biliknya}
berepolarisasi.
\end{solution}

\begin{solution}{exo:b2:heart:4}
Di dalam jangkauan kerjanya \omterm{def:b2:heart:anatomy}{biliknya} menyemburkan lebih banyak ketika
ia telah diisi lebih banyak. Jika \omterm{def:b2:heart:anatomy}{bilik} kanannya memompa lebih banyak,
yang kiri menerima lebih banyak beberapa denyut kemudian lalu, dengan
hukum yang sama, memompa lebih banyak: ketakseimbangan apa pun
membetulkan dirinya sendiri tanpa sebuah isyarat.
\end{solution}

\begin{solution}{exo:b2:heart:5}
\omterm{prop:b2:heart:cycle}{Volume sekuncupnya} \qty{80}{mL}; \omterm{prop:b2:heart:cycle}{fraksi semburannya} $80/130 = \qty{62}{\%}$;
keluarannya $80\times 70 = \qty{5.6}{L/min}$. Yang gagal: \qty{60}{mL},
$60/180 = \qty{33}{\%}$, \qty{4.2}{L/min} --- \omterm{def:b2:heart:anatomy}{bilik} yang melebar itu
menjaga keluarannya tetap hampir cukup dengan bekerja pada volume yang
besar, tetapi hanya menyemburkan sepertiga dari yang ia tahan.
\end{solution}

\begin{solution}{exo:b2:heart:6}
$W = 100\times 133\times 7\times 10^{-5} = \qty{0.93}{J}$; dayanya
$0.93\times 1.2 = \qty{1.1}{W}$. Pada \qty{150}{mmHg}: \qty{1.4}{J},
\qty{1.7}{W}. Tambahan \qty{0.56}{W} mekanis adalah \qty{3.7}{W}
metabolis, yaitu \qty{0.19}{mL} oksigen tiap detik: \qty{11}{mL/min} lebih banyak.
\end{solution}

\begin{solution}{exo:b2:heart:7}
$60/0.5 = 120$ denyut tiap menit; $60/1.5 = 40$. Yang ketiga: blok
\omterm{def:b2:heart:anatomy}{jantung} lengkap --- \omterm{def:b2:heart:anatomy}{serambinya} berdenyut pada 100 di bawah nodusnya,
\omterm{def:b2:heart:anatomy}{biliknya} pada 33 di bawah pacunya sendiri, dan nodus AV-nya tidak
menghantarkan apa pun.
\end{solution}

\begin{solution}{exo:b2:heart:8}
Hanyutannya $20/25 = \qty{0.8}{s}$, ditambah \qty{0.15}{s}:
\qty{0.95}{s}, yaitu 63 denyut tiap menit. Asetilkolin: $25/18 = \qty{1.39}{s}$, ditambah $0.15$:
\qty{1.54}{s}, yaitu 39 tiap menit. Noradrenalin: $20/40 = \qty{0.5}{s}$,
ditambah $0.15$: \qty{0.65}{s}, yaitu 92 tiap menit.
\end{solution}

\begin{solution}{exo:b2:heart:9}
Tundaannya membiarkan \omterm{def:b2:heart:anatomy}{serambinya} merampungkan pengosongan ke dalam
\omterm{def:b2:heart:anatomy}{biliknya} sebelum \omterm{def:b2:heart:anatomy}{biliknya} berkerut; tanpa tundaan itu \omterm{def:b2:heart:anatomy}{serambi} dan
\omterm{def:b2:heart:anatomy}{biliknya} akan berkerut bersama-sama, katup AV-nya akan menutup pada
\omterm{def:b2:heart:anatomy}{bilik} yang terisi separuh, dan \omterm{prop:b2:heart:cycle}{volume sekuncupnya} akan turun sebesar
sumbangan \omterm{def:b2:heart:anatomy}{serambinya}. Selang PR yang panjang hanya menunda sistol
\omterm{def:b2:heart:anatomy}{biliknya} dan tidak berongkos apa pun sampai tundaannya tumbuh begitu
panjang sehingga denyutnya hilang atau penghantarannya gagal sama
sekali.
\end{solution}

\begin{solution}{exo:b2:heart:10}
Istirahat: $5000/45 = \qty{111}{mL}$; olahraga: $\num{30000}/180 =
\qty{167}{mL}$. Hukum Starlingnya mengubah kembalian \omterm{def:b2:blood-circulation:vessels}{vena} yang lebih
besar ketika berolahraga menjadi \omterm{prop:b2:heart:cycle}{volume sekuncup} yang lebih besar;
saraf simpatisnya menambahkan laju dan daya kerut (volume akhir sistol
yang lebih rendah). Di atas kira-kira 200 diastolnya terlalu pendek
untuk mengisi \omterm{def:b2:heart:anatomy}{biliknya} dan \omterm{prop:b2:heart:cycle}{volume sekuncupnya} turun lebih cepat
daripada kenaikan lajunya.
\end{solution}

\begin{solution}{exo:b2:heart:11}
$v = Q/A$: $250/3 = \qty{83}{cm/s}$ dan $250/1 = \qty{250}{cm/s}$.
$\Delta P = \tfrac12\times 1060\times v^{2}$: \qty{365}{Pa}
(\qty{2.7}{mmHg}) dan \qty{3300}{Pa} (\qty{25}{mmHg}). \omterm{def:b2:heart:anatomy}{Biliknya}
membangkitkan tekanan tambahannya dengan menebalkan dindingnya; \omterm{def:b2:heart:anatomy}{bilik}
yang tebal dan kaku itu akhirnya melampaui pasokan koronernya dan
terisi dengan buruk, lalu gagal.
\end{solution}

\begin{solution}{exo:b2:heart:12}
Keotomatisannya memberikan sebuah denyut tanpa masukan apa pun; hukum
Starlingnya menyerasikan keluarannya dengan kembaliannya dan
mengimbangkan kedua \omterm{def:b2:heart:anatomy}{biliknya} tanpa masukan apa pun; saraf dan
adrenalinnya hanya menetapkan laju dan daya kerut di sekitar inti yang
mengatur dirinya sendiri itu. Sebuah \omterm{def:b2:heart:anatomy}{jantung} yang dicangkokkan
berdenyut (pada kira-kira 100, tanpa tonus vagus), menyetel \omterm{prop:b2:heart:cycle}{volume
sekuncupnya} kepada pengisiannya, dan dapat menaikkan keluarannya ketika
berolahraga lewat mekanisme Starlingnya dan lewat adrenalin yang
beredar --- tetapi perlahan, dan kurang daripada \omterm{def:b2:heart:anatomy}{jantung} yang bersaraf.
\end{solution}

\begin{solution}{pb:b2:heart:1}
\textbf{1.} \qty{70}{mL}; $70/130 = \qty{54}{\%}$; $70\times 70 =
\qty{4.9}{L/min}$.
\textbf{2.} $60/70 = \qty{0.86}{s}$; diastolnya $0.86 - 0.3 =
\qty{0.56}{s}$.
\textbf{3.} $70/0.3 = \qty{233}{mL/s}$; $233/3 = \qty{78}{cm/s}$.
\textbf{4.} $70/1500 = \qty{0.047}{s}$, kira-kira \qty{50}{ms}.
\textbf{5.} $100\times 133\times 7\times 10^{-5} = \qty{0.93}{J}$;
$0.93\times 70/60 = \qty{1.1}{W}$.
\textbf{6.} Kedua \omterm{def:b2:heart:anatomy}{biliknya} $1.1\times 1.2 = \qty{1.3}{W}$; pada
\qty{15}{\%}, \qty{8.7}{W} metabolis; $8.7/20 = \qty{0.43}{mL}$
oksigen tiap detik, \qty{26}{mL/min}.
\textbf{7.} Dengan menyaripatikan \qty{0.14}{mL} tiap mililiter \omterm{def:b2:blood-circulation:blood}{darah}:
$26/0.14 = \qty{190}{mL/min}$.
\textbf{8.} Selangnya \qty{0.86}{s}, waktu hanyutannya \qty{0.71}{s}:
$20/0.71 = \qty{28}{mV/s}$.
\textbf{9.} Selangnya \qty{0.6}{s}, hanyutannya \qty{0.45}{s}: \qty{44}{mV/s}.
\textbf{10.} Selangnya \qty{0.33}{s}, hanyutannya \qty{0.18}{s}: \qty{110}{mV/s}.
\textbf{11.} R--R \qty{0.86}{s}; P: depolarisasi \omterm{def:b2:heart:anatomy}{serambi}; QRS:
depolarisasi \omterm{def:b2:heart:anatomy}{bilik}; T: repolarisasi \omterm{def:b2:heart:anatomy}{bilik}.
\textbf{12.} Tundaannya membiarkan \omterm{def:b2:heart:anatomy}{serambinya} mengosongkan diri ke
dalam \omterm{def:b2:heart:anatomy}{biliknya} sebelum \omterm{def:b2:heart:anatomy}{biliknya} berkerut; pada 180 seluruh daurnya
adalah \qty{0.33}{s}, sehingga tundaan \qty{0.16}{s} yang tak berubah
akan menghabiskan separuhnya --- saraf simpatisnya mempercepat
penghantaran nodusnya juga.
\textbf{13.} Blok \omterm{def:b2:heart:anatomy}{jantung} lengkap: \omterm{def:b2:heart:anatomy}{biliknya} pada 40 tiap menit, dipacu
oleh berkas atau \omterm{prop:b2:heart:conduction}{serabut Purkinjenya}, sedangkan \omterm{def:b2:heart:anatomy}{serambinya} pada 75 di
bawah \omterm{prop:b2:heart:conduction}{nodus sinoatriumnya}.
\textbf{14.} \qty{120}{mL}; $120/150 = \qty{80}{\%}$; $120\times 180 =
\qty{21.6}{L/min}$.
\textbf{15.} $0.33 - 0.3 = \qty{0.03}{s}$: nyaris tanpa waktu untuk
terisi; lebih cepat lagi dan \omterm{prop:b2:heart:cycle}{volume sekuncupnya} runtuh.
\textbf{16.} $W = 120\times 133\times 1.2\times 10^{-4} = \qty{1.9}{J}$,
$\times 3$ tiap detik: \qty{5.7}{W}; kedua \omterm{def:b2:heart:anatomy}{biliknya} \qty{6.9}{W};
metabolisnya \qty{46}{W}; oksigennya \qty{2.3}{mL/s}, \qty{140}{mL/min}.
\textbf{17.} $140/0.14 = \qty{1000}{mL/min}$, lima kali aliran ketika
beristirahat, yang harus diantarkan dalam diastol yang menyusut menjadi
sepersepuluh daurnya: \omterm{def:b2:blood-circulation:vessels}{arteriol} koronernya harus melebar sampai batasnya.
\textbf{18.} Lajunya kira-kira 100 tanpa tonus vagus, volume akhir
sistolnya tetap \qty{60}{mL}: \omterm{prop:b2:heart:cycle}{volume sekuncupnya} $150 - 60 = \qty{90}{mL}$,
keluarannya \qty{9}{L/min}.
\textbf{19.} Lajunya sendiri ($70 \to 180$ pada \qty{70}{mL}, yakni
\qty{12.6}{L/min}): $+7.7$; pengisiannya (\qty{20}{mL} lebih banyak pada 180):
$+3.6$; daya kerutnya (\qty{30}{mL} lebih sedikit tersisa pada 180): $+5.4$ ---
dari 4.9 menjadi \qty{21.6}{L/min}, ketiga sumbangannya berjumlah
tepat $+16.7$.
\textbf{20.} $233/1 = \qty{233}{cm/s}$, \qty{2.3}{m/s}.
\textbf{21.} $\tfrac12\times 1060\times 2.33^{2} = \qty{2900}{Pa}$,
\qty{22}{mmHg}.
\textbf{22.} $120/0.3 = \qty{400}{mL/s}$, \qty{4}{m/s}; $\tfrac12\times
1060\times 16 = \qty{8500}{Pa}$, \qty{64}{mmHg}.
\textbf{23.} Puncaknya $120 + 22 = \qty{142}{mmHg}$ ketika
beristirahat, kira-kira $140 +
64 = \qty{204}{mmHg}$ ketika
berolahraga; tekanan semburan reratanya ketika beristirahat naik dari
100 menjadi 122: kerja tiap denyutnya $\times 1.22$.
\textbf{24.} Dinding yang lebih tebal membangkitkan gaya yang lebih
besar dan menurunkan tegasan tiap serabutnya; tetapi bongkah yang harus
dialiri bertambah sementara aliran koronernya, yang diperas dinding
tebalnya pada sistolnya dan diberi diastol yang lebih pendek, tidak
mengimbanginya, dan dinding yang tebal mengendur dengan buruk lalu
terisi lebih sedikit: iskemia dan gagal \omterm{def:b2:heart:anatomy}{jantung} menyusul.
\textbf{25.} \qty{4.9}{L/min} ketika beristirahat, \qty{21.6}{L/min}
ketika berolahraga; \qty{0.93}{J} tiap denyut; landaiannya
\qty{22}{mmHg} ketika beristirahat dan \qty{64}{mmHg} ketika
berolahraga.
\end{solution}
